%======================================================================
%  The node: Lenovo ThinkSystem SD650 V3
%======================================================================

\section{The SD650 V3 node}
\label{sec:sd650-node}

The processors are housed in Lenovo ThinkSystem SD650~V3 trays.  Each
tray carries two independent dual-socket nodes, six trays share a 6U
DW612S enclosure, and the unit that Slurm allocates is one
node~\cite{lp1603}.  \figref{fig:sd650-internal} shows a tray.  The HBM
is not visible in it: it sits inside the processor package under the
cold plate, which is why its capacity is fixed per socket.

\begin{figure}[htbp]
\centering
% Crop the border and drop shadow embedded in the source image.
\includegraphics[width=0.94\textwidth,trim=6bp 27bp 27bp 6bp,clip]{sd650v3-internal-view.png}
\caption{An SD650~V3 tray: two dual-socket nodes, each processor flanked
by eight DIMM slots, all cooled by water through copper cold plates.
Reproduced from the Lenovo ThinkSystem SD650~V3 product
guide~\cite{lp1603}.}
\label{fig:sd650-internal}
\end{figure}

The node is cooled by direct water cooling, without fans.  This supports
sustained all-core load in a dense enclosure, but it does not fix the
clock frequency: the package power limit and the frequency reductions
for vector and matrix code still apply (\secref{sec:xeon-max-peaks}).
Frequency is therefore recorded during every measurement rather than
assumed.

\paragraph{Power telemetry.}  The SD650~V3 has an XClarity management
controller and a TI INA226-based power meter that can sample at
100\,Hz~\cite{lp1603}.  On CRESCO8 these are administrative interfaces.
What a user can actually read was established on an allocated node:

\begin{itemize}
  \item \textbf{Slurm energy accounting} uses the IPMI plugin with a
        60-second sampling interval.  It reports whole-node energy for a
        job and the node's current and average power; an idle HBM node
        reported about 390\,W.
  \item \textbf{RAPL} package and DRAM energy counters exist but are
        readable only by root: reading \texttt{energy\_uj} under
        \texttt{/sys/class/powercap} fails with a permission error.
  \item \textbf{Model-specific registers} and therefore
        \texttt{turbostat} are likewise restricted to root.
\end{itemize}

The only energy measurement available to the experiments is consequently
whole-node energy at one-minute granularity.  It cannot separate
processor from memory power, nor attribute energy to an inference phase
shorter than a minute.  Energy per token is therefore outside the scope
of the measurements unless the administrators grant RAPL access, and
this chapter makes no energy claims.
